\subsection{Reloj: PIT + PIC $\to$ MTIP + SSIP}
\label{sec:reloj}

\subsubsection*{Qué hacía el TP x86}

La \texttt{\_isr32} llamaba a \texttt{nextClock} (animación del cursor) y \texttt{pic\_finish1} (EOI al PIC maestro), y
ejecutaba \texttt{iret}.

\subsubsection*{Equivalente en RISC-V}

En RISC-V el timer está en CLINT/ACLINT~\cite{riscv-aclint}: dos MMIOs, \texttt{mtime} (contador libre) y
\texttt{mtimecmp} (comparador por HART). Cuando \texttt{mtime} $\geq$ \texttt{mtimecmp}, se levanta \textbf{MTIP}
(machine timer interrupt pending). El problema es que MTIP \textbf{solo} se puede atender en M-mode (no se puede delegar
a S-mode directamente~\cite{riscv-priv}). Un sistema con OpenSBI resolvería esto usando
\texttt{sbi\_set\_timer}~\cite{riscv-sbi}; con \texttt{-bios none} no hay SBI y debemos resolverlo nosotros.

El patrón estándar (el mismo que usa xv6 en su build sin SBI~\cite{xv6-riscv} y que documenta el wiki OSDev para
RISC-V~\cite{osdev-riscv}) es:

\begin{enumerate}
  \item Atender el MTIP en M-mode (\texttt{machine\_trap\_entry}, \texttt{mcause}$=$7).
  \item Reprogramar el próximo \texttt{mtimecmp}.
  \item \textbf{Levantar un supervisor software interrupt (SSIP)} escribiendo \texttt{sip.SSIP}$=$1 desde M-mode.
  \item Ejecutar \texttt{mret} y dejar que la CPU dispare el handler S-mode (vía \texttt{stvec}) con \texttt{scause}$=$1
  (SSIP).
\end{enumerate}

En \texttt{machine.S}:

\begin{codesnippet}
\begin{verbatim}
 # Reprogramar mtimecmp
 csrr t0, time
 la t1, g_timer_delta_ticks
 ld t1, 0(t1)
 add t0, t0, t1
 csrr t1, mhartid
 slli t1, t1, 3
 li t2, MTIMECMP_BASE        # 0x0200_4000 + 8*hartid
 add t2, t2, t1
 sd t0, 0(t2)

 # Levantar SSIP
 li t0, 2
 csrs sip, t0
\end{verbatim}
\end{codesnippet}

Y del lado de S-mode, \texttt{trap\_handle} reconoce SSIP como el tick:

\begin{codesnippet}
\begin{verbatim}
 if (is_interrupt) {
     if (cause_code == 1) {  // SSIP
         const uint64_t ssip_mask = (1ull << 1);
         asm volatile("csrc sip, %0" : : "r"(ssip_mask));
         /* cooldowns + poll UART + spinner + scheduler */
         return schedule_next();
     }
     return tf;
 }
\end{verbatim}
\end{codesnippet}

El \texttt{csrc sip, SSIP} es el análogo de \texttt{pic\_finish1}: limpia el bit pendiente para que el interrupt no se
vuelva a disparar. La función \texttt{nextClock} del TP original (animación del spinner en la esquina) se preserva con
\texttt{update\_spinner()} en \texttt{kernel.cpp}.

\subsection{Teclado: IRQ1 + scancodes $\to$ UART polling}

\subsubsection*{Qué hacía el TP x86}

La \texttt{\_isr33} leía el scancode de \texttt{in al, 0x60}, pusheaba \texttt{eax} y llamaba a
\texttt{idt\_int\_teclado} (que traducía el scancode a la acción del juego), luego \texttt{pic\_finish1} y
\texttt{iret}.

\subsubsection*{Equivalente en RISC-V}

QEMU \texttt{virt} no tiene teclado PS/2 estándar; lo que hay es el UART 16550 sobre \texttt{-serial stdio}. Podría
tomarse la interrupción del UART vía PLIC, pero para no sumar un mecanismo más al informe terminé haciéndolo por
\textbf{polling en cada tick}, que a la frecuencia del reloj ($\sim$20~Hz) se percibe igual. En lugar de un ISR por
tecla, \texttt{trap\_handle} llama a \texttt{poll\_uart\_input()} al atender el tick, y ese drena el buffer con
\texttt{uart::getc\_nonblock()}:

\begin{codesnippet}
\begin{verbatim}
 int uart::getc_nonblock() {
     if ((mmio_read8(UART_LSR) & LSR_DATA_READY) == 0)
         return -1;
     return mmio_read8(UART_RHR);
 }
\end{verbatim}
\end{codesnippet}

Al llegar un byte se lo pasa a \texttt{handle\_key}, que mapea teclas (no scancodes) a acciones de juego:

\begin{itemize}
  \item Jugador ROJO: \texttt{W} (arriba), \texttt{S} (abajo), \texttt{Q}/\texttt{E} (spawnear minero).
  \item Jugador AZUL: \texttt{I} (arriba), \texttt{K} (abajo), \texttt{O}/\texttt{P} (spawnear minero).
\end{itemize}

El byte recibido se imprime en la esquina superior derecha en hex (equivalente al ``mostrar la tecla 0..9 en pantalla''
del TP x86), útil para validar input.

\subsection{\texttt{int 47} (x86) $\to$ \texttt{ecall} (RISC-V)}

En x86 se declaraba una Interrupt Gate con DPL=3 en la entrada 47 de la IDT, y la \texttt{\_isr47} inicialmente seteaba
\texttt{eax}$=$\texttt{0x42} y hacía \texttt{iret}. En RISC-V el análogo de esa \emph{software interrupt} es
\texttt{ecall}: desde U-mode genera un trap con \texttt{scause}$=$8 (``Environment call from U-mode''). No hay un
``número'' configurable como en \texttt{int N}; lo que sí se mantiene del TP3 son los IDs de servicio, que en este port
van por \texttt{a7} siguiendo la convención estándar de syscalls de RISC-V:

\begin{center}
\begin{tabular}{l|l}
\textbf{Servicio} & \textbf{ABI} \\
\hline
\texttt{move(d)} & \texttt{a7}=177788, \texttt{a0}=\texttt{d} \\
\texttt{take()}  & \texttt{a7}=310311, retorna en \texttt{a0} \\
\texttt{getId()} & \texttt{a7}=760279, retorna en \texttt{a0} \\
\end{tabular}
\end{center}

No se usa un ``eax transitorio con \texttt{0x42}'' porque el handler puede escribir directamente sobre \texttt{tf->a0} y
el \texttt{sret} restaura \texttt{a0} como valor de retorno. El stub de usuario vive en \texttt{riscv/syscall.h}:

\begin{codesnippet}
\begin{verbatim}
 static inline int32_t syscall_take() {
     register uint64_t a7 asm("a7") = 310311;
     register uint64_t a0 asm("a0");
     asm volatile("ecall" : "=r"(a0) : "r"(a7) : "memory");
     return (int32_t)a0;
 }
\end{verbatim}
\end{codesnippet}

El despacho en \texttt{trap\_handle} se describe en el ejercicio~7.
